\documentclass[11pt,oneside]{memoir}
\usepackage{raeez-math-template}
\title{Ternary comparison on the collision diagonals}
\author{}
\date{}
\begin{document}
\maketitle
\begin{abstract}
The nonadjacent collision of three currents requires a correction on a partial diagonal.
An explicit contraction of the Heisenberg Chevalley--Cousin complex supplies this correction at nonzero level.
The resulting cochain map includes all six ordered three-letter words and all twelve mixed words.
It retains the prescribed mixed intermediate images and every normal derivative.
Weighted edge integrals extend the polynomial comparison to every global ternary coefficient and every finite inner pole order.
Induction gives a comparison of right differential modules with a specified source and connection.
The target at nonzero level is contractible, so the construction does not identify the original Heisenberg state cohomology.
\end{abstract}
\providecommand{\HH}{\mathcal H}
\providecommand{\pp}{\operatorname{pp}}
\chapter{Ternary comparison on the collision diagonals}
\label{chap:bd26-comparison}

\section{The nonadjacent current collision}

Three ordered letters have two adjacent multiplication faces.
Three chiral insertions have three pair-collision maps.
Their comparison must account for the collision between the first and third points.
For Heisenberg currents, its coefficient is the level times a normal derivative of the vacuum delta class.
A primitive on that partial diagonal supplies the missing term.

Work over $S=\mathbb C[\kappa]$ on $X=\mathbb A^1_{\mathbb C}$.
The parameter $\kappa$ and all position coordinates have degree zero.
All tensor products use $S$ unless another base is indicated.
The shift convention is $M[r]^n=M^{n+r}$ with differential $(-1)^r d_M$.
Parity is cohomological degree modulo two.
No state-space or parameter completion is used.

The Heisenberg state space, vacuum, current, and translation are
\[
 V=S[x_1,x_2,\ldots],\qquad \mathbf1=1,\qquad J=x_1,
 \qquad T=\sum_{n\geq1}n x_{n+1}\partial_{x_n}.
\]
Its states are even and have degree zero.
Every state uses finitely many variables.
The current modes act by
\[
 J_{-n}=x_n\cdot,\qquad J_n=\kappa n\partial_{x_n}\quad(n\geq1),
 \qquad J_0=0.
\]
Put $J(z)=\sum_{m\in\mathbb Z}J_mz^{-m-1}$.
The field of $x_{n_1}\cdots x_{n_r}$ is the normally ordered product of
$\partial_z^{n_i-1}J(z)/(n_i-1)!$.
Normal ordering places every multiplication operator before every derivative operator.
Write this field as $Y(a,z)$.

Commuting a derivative past a multiplication operator gives the contraction
\[
 C_{nm}(t)=
 \frac{\kappa(-1)^{n-1}(n+m-1)!}{(n-1)!(m-1)!}\,t^{-n-m}.
\]
For any finite set of input monomials, the field product is a finite sum over partial pairings.
Each pair contributes $C_{nm}$, and unpaired factors remain normally ordered.
This follows by applying $AB=BA+[A,B]$ at each crossing.
A contracted factor is removed and cannot occur in another crossing.
The formula gives $Y(a,t)b\in V((t))$ with a finite lower bound.
In particular,
\begin{equation}\label{eq:bd26-field}
 p(t):=Y(J,t)J=\frac\kappa{t^2}+(e^{tT}J)J,
 \qquad e^{tT}J=\sum_{r\geq0}x_{r+1}t^r.
\end{equation}

Introduce the differential graded vertex superalgebra
\begin{equation}\label{eq:bd26-state}
 W=V\otimes_S\Lambda_S(\epsilon),\qquad
 |\epsilon|=-1,\qquad d_W\epsilon=\kappa,\qquad T\epsilon=0.
\end{equation}
For homogeneous scalar factors $r,s\in\Lambda_S(\epsilon)$, set
$Y_W(ar,t)(bs)=Y(a,t)b\,rs$.
The differential is a graded derivation and commutes with translation.
The scalar $\epsilon$ is odd and squares to zero.

The associated right differential module is $\HH_W=\omega_X\otimes W$, with
\[
 (fa\,dz)\partial_z=-(f'a+fTa)\,dz.
\]
For $t=z-w$ and $y=w$, the chiral operation is
\begin{equation}\label{eq:bd26-chiral}
 \mu_W\bigl(f(z,w)(a\,dz\boxtimes b\,dw)\bigr)
 =\pp_t\bigl(f(y+t,y)Y_W(a,t)b\bigr).
\end{equation}
The input coefficient may have finite poles at $t=0$.
The operation multiplies the Laurent series before taking its finite principal part.
The normal convention and orientation are
\begin{equation}\label{eq:bd26-normal}
 \pp_t\frac{a(y)}{t^{r+1}}
 =\frac{a(y)\,dy}{r!}\delta_\Delta\partial_t^r,
 \qquad dz\wedge dw=dt\wedge dy.
\end{equation}
There is no multiplication by $t^{-1}$ on an existing principal-part class.
Regular position functions act before passage to the quotient.

For clarity, the chiral identities can be checked directly on these formulas.
Translation gives $Y(Ta,t)=\partial_tY(a,t)$ and
$TY(a,t)b=\partial_tY(a,t)b+Y(a,t)Tb$.
They identify both input right derivatives with the corresponding output right derivatives in \eqref{eq:bd26-chiral}.
The contraction satisfies $C_{mn}(-t)=C_{nm}(t)$.
Interchanging the two Wick inputs and translating the result by $e^{tT}$ therefore gives graded vertex skew symmetry.
On normal principal parts, coordinate transposition acts by
$p(t,y)\mapsto-\pp_t(e^{tT}p(-t,y+t))$.
This follows on $a/t$ and $a/t^2$ from \eqref{eq:bd26-normal}, and then on all normal derivatives from the Weyl relations.
Together these formulas prove graded skew symmetry of $\mu_W$.
For three monomials, the Wick sum belongs to
\[
 W[[u,v]][u^{-1},v^{-1},(u-v)^{-1}].
\]
Its two radial expansions give the two ordered products.
Pairing the first two inputs first gives its expansion at $u=v$, and hence the iterate.
Partial fractions in $u$, followed by the residue in $v$, give
\[
 \operatorname{Res}_v\operatorname{Res}_t G(v+t,v)\,dt\,dv
 =\operatorname{Res}_u\operatorname{Res}_v i_{u,v}G(u,v)\,dv\,du
 -\operatorname{Res}_v\operatorname{Res}_u i_{v,u}G(u,v)\,du\,dv.
\]
Here $i_{u,v}$ expands in nonnegative powers of $v/u$.
For a power-series numerator, each residue depends on finitely many coefficients.
Apply the identity after multiplication by every polynomial testing the normal derivatives.
It proves chiral Jacobi as an equality in the small-diagonal transfer module.
Graded commutativity of the scalar factors supplies the super signs.
Finally, $Y_W(\mathbf1,t)a=a$ gives the normalized unit operation $\pp_t(fa)$.

\section{The five collision summands}

Let $I=\{1,2,3\}$ and $R=S[z_1,z_2,z_3]$.
For a partition $\pi$ of $I$, let $\Delta_\pi:X^\pi\hookrightarrow X^3$ identify positions within each block.
Let $j_\pi$ include the locus of distinct block positions in $X^\pi$.
The symbol $(\Delta_\pi)_!$ denotes closed direct image for right differential modules, using the transfer module.
It includes all finite normal derivatives and no additional cohomological shift.
Define
\begin{equation}\label{eq:bd26-target}
 E=
 \bigoplus_{\pi\vdash I}
 (\Delta_\pi)_!(j_\pi)_*j_\pi^*
 \left(\boxtimes_{B\in\pi}\HH_W[1]\right).
\end{equation}
This is the three-point Chevalley--Cousin module
\cite{BDdiagonal}*{\S3.4.11, equation (3.4.11.1)}.
On the affine space $X^3$, the same notation denotes its module of global sections.
The module formulas determine the associated quasi-coherent sheaf maps.
The summands correspond to $1|2|3$, $12|3$, $13|2$, $1|23$, and $123$.
The block order is auxiliary.
Interchanging blocks uses the Koszul sign of their suspended state degrees.

Write $[a_1|\cdots|a_r]$ for suspended states, with
$|sa|=|a|-1$ and $d(sa)=-s(d_Wa)$.
The total differential is $Q=Q_1+Q_2$.
The operator $Q_1$ is the tensor differential of the suspended states.
The binary collision has sign
\[
 Q_2[a|b]=(-1)^{|a|}[\mu_W(a,b)].
\]
For three blocks, with $\alpha=|a|$, $\beta=|b|$, and $\gamma=|c|$, it is
\begin{equation}\label{eq:bd26-collision-signs}
 \begin{aligned}
 Q_2[a|b|c]={}&(-1)^\alpha[\mu_{12}(a,b)|c]\\
 &+(-1)^{(\beta-1)(\gamma-1)+\alpha}[\mu_{13}(a,c)|b]\\
 &+(-1)^{(\alpha-1)(\beta+\gamma-2)+\beta}[\mu_{23}(b,c)|a].
 \end{aligned}
\end{equation}
The three maps first localize away from collisions with the remaining point.
They then apply \eqref{eq:bd26-chiral} and the appropriate direct image.
Every coordinate permutation includes its density identification.
Thus \eqref{eq:bd26-collision-signs} uses the full common three-point source.

The differential Leibniz identity gives $Q_1Q_2+Q_2Q_1=0$.
Applying $Q_2$ twice to \eqref{eq:bd26-collision-signs} gives
\[
 (-1)^\beta[
 \mu(\mu(a,b),c)-\mu(a,\mu(b,c))
       +(-1)^{\alpha\beta}\mu(b,\mu(a,c))].
\]
This vanishes by the chiral Jacobi identity just proved.
Hence $Q^2=0$, including every normal derivative.
Both $Q_1$ and $Q_2$ commute with the right differential actions.

For the small diagonal, put $t=z_1-z_2$, $s=z_2-z_3$, and $y=z_3$.
Its unsuspended principal parts are
\[
 P_3(W)=W\otimes_S
 \frac{S[y,t,s,t^{-1},s^{-1}]}
 {S[y,t,s,t^{-1}]+S[y,t,s,s^{-1}]}.
\]
The orientation and finite normal basis are
\begin{equation}\label{eq:bd26-small}
 \begin{gathered}
 dz_1\wedge dz_2\wedge dz_3=dt\wedge ds\wedge dy,\\
 \pp\frac{a(y)}{t^{r+1}s^{l+1}}
 =\frac{a(y)\,dy}{r!l!}\delta_3\partial_t^r\partial_s^l.
 \end{gathered}
\end{equation}
Its right actions are $-\partial_t$, $-\partial_s$, and $-\partial_y-T$.
The small-diagonal summand of $E$ is $P_3(W)[1]$.
Its differential is $-d_W$, and its only degrees are $-2$ and $-1$.

\section{A contraction preserving each diagonal}

Distinguish label $1$.
For each partition, the block containing $1$ determines an odd scalar action.
If that block occupies position $j$ in the displayed block order, define
\begin{equation}\label{eq:bd26-epsilon-operator}
 \mathsf e[a_1|\cdots|a_r]
 =(-1)^{\sum_{i<j}(|a_i|-1)}
 [a_1|\cdots|\epsilon a_j|\cdots|a_r].
\end{equation}
The formula extends to localized coefficients and normal derivatives.
It has degree $-1$.
Changing the block order gives the same operator because \eqref{eq:bd26-epsilon-operator} uses the tensor sign for an odd map.

\begin{lemma}\label{lem:bd26-scalar-homotopy}
On the entire five-summand complex,
\begin{equation}\label{eq:bd26-homotopy-level}
 \mathsf e^2=0,\qquad
 Q_1\mathsf e+\mathsf e Q_1=-\kappa\,\mathrm{id},\qquad
 Q_2\mathsf e+\mathsf e Q_2=0.
\end{equation}
The operator $\mathsf e$ commutes with the right differential actions and preserves every partition summand.
Consequently, over $S_*=S[\kappa^{-1}]$,
\begin{equation}\label{eq:bd26-contraction}
 h=-\kappa^{-1}\mathsf e,
 \qquad h^2=0,\qquad Qh+hQ=\mathrm{id}_E.
\end{equation}
\end{lemma}
\begin{proof}
The first identity follows from $\epsilon^2=0$.
On one suspended factor, $Q_1\mathsf e+\mathsf e Q_1=-\kappa$ follows from
$d_W(\epsilon a)=\kappa a-\epsilon d_Wa$ and the suspension minus sign.
The tensor signs cancel the differentials on the other factors.
This proves the second identity in every summand.

The chiral operation is graded linear in the scalar factors:
\[
 \mu_W(\epsilon a,b)=\epsilon\mu_W(a,b),\qquad
 \mu_W(a,\epsilon b)=(-1)^{|a|}\epsilon\mu_W(a,b).
\]
For two blocks with the distinguished block first, applying $Q_2\mathsf e$ has coefficient $(-1)^{|a|-1}$.
Applying $\mathsf e Q_2$ has coefficient $(-1)^{|a|}$.
They cancel.
With the distinguished block second, the tensor sign and the second displayed scalar identity give the same cancellation.

For three blocks, order the distinguished block first and use \eqref{eq:bd26-collision-signs}.
For pairs $12$ and $13$, replacing $a$ by $\epsilon a$ changes the displayed sign by $-1$.
The distinguished block remains first in each output, so the two compositions cancel.
For pair $23$, put $N=(\alpha-1)(\beta+\gamma-2)+\beta$.
The $Q_2\mathsf e$ coefficient is $(-1)^{N+\beta+\gamma}$.
In the output $[\mu(b,c)|a]$, the $\mathsf e Q_2$ coefficient is
$(-1)^{N+\beta+\gamma-1}$.
They also cancel.
This checks every collision in $E$.

The scalar is constant in position and translation fixes it.
It therefore commutes with tangent and normal right derivatives, and with every regular coordinate multiplication.
It changes no partition, proving the support assertion.
Division by $-\kappa$ proves \eqref{eq:bd26-contraction}.
\end{proof}

The homotopy retains all partial diagonals.
Its restriction to the small diagonal sends $[a]$ to $-\kappa^{-1}[\epsilon a]$.
On a partial diagonal it inserts $\epsilon$ in the block containing label $1$.
Thus it can cancel a nonadjacent collision before reaching the small diagonal.

\section{The ordered three-input complex}

Let $\mathcal O_I$ be the six total orders of $I$.
For nonempty $A\subset\mathcal O_I$, let $R_A$ invert $z_i-z_j$ exactly when the relative order of $i,j$ varies in $A$.
Let
\[
 \Omega_A=S[u_\sigma,du_\sigma\mid\sigma\in A]
 /(\textstyle\sum u_\sigma-1,\sum du_\sigma),
\]
where the $du_\sigma$ generate an exterior algebra.
Define $C$ to be the algebra of families in $R_A\otimes_S\Omega_A$ compatible under all face restrictions and coefficient localizations.
Its differential $d_C$ fixes positions and sends $u_\sigma$ to $du_\sigma$.
Put
\[
 q_{ij}=\sum_{i\prec_\sigma j}u_\sigma,
 \qquad \omega_{ij}=\frac{dq_{ij}}{z_i-z_j}.
\]
On a component without the denominator, the numerator vanishes and this component is zero.
These formulas define closed compatible families.

The coefficient-weighted state algebra has generators $a_i$ of degree zero and $b_{ij}$ of degree $-1$.
The states are closed, and coefficients act graded centrally.
Its nonzero nonunit products are $a_i a_j=\omega_{ij}b_{ij}$ for $i\ne j$.
Every product involving a $b$ and another nonunit state vanishes, and $a_i^2=0$.
The labels $b_{ij}$ and $b_{ji}$ are distinct.
These products are associative because every product of three nonunit states vanishes.

For distinct $i,j,k$, the multilinear three-input bar sector contains
\[
 e_{ijk}=[a_i|a_j|a_k],\qquad
 m_{ij|k}=[b_{ij}|a_k],\qquad
 n_{k|ij}=[a_k|b_{ij}].
\]
There are six $e$-words and twelve mixed words.
Each displayed word has degree $-3$ before its coefficient is included.
Let $O$ be their direct sum of free $C$-modules.
The differential is
\begin{equation}\label{eq:bd26-source}
 \begin{gathered}
 be_{ijk}=-\omega_{ij}m_{ij|k}-\omega_{jk}n_{i|jk},
 \qquad bm_{ij|k}=bn_{k|ij}=0,\\
 b(cv)=d_Cc\,v+(-1)^{|c|}c\,bv.
 \end{gathered}
\end{equation}
These signs follow from $s(ca)=(-1)^{|c|}c\,sa$ and
$b_2(sa\otimes sb)=(-1)^{|a|}s(ab)$.
Closedness of the $\omega_{ij}$ gives $b^2=0$ directly on every summand.
Thus $O$ is the actual multilinear sector of the relative ordered bar.

For $m_{ij|k}$, restrict the compatible family to the edge $\{ijk,jik\}$.
For $n_{k|ij}$, use the edge $\{kij,kji\}$.
Set $u=q_{ij}$ and $t=z_i-z_j$ on either edge.
The edge ring is exactly $R[t^{-1}]$.
In degree zero, the restriction is $f(u)$ with $f(0),f(1)\in R$.
In degree one, it is $g(u)du$ with $g\in R[t^{-1}][u]$.
Higher-degree restrictions vanish.
These are restrictions of entire families, and each restriction commutes with $d_C$.

Write
\begin{equation}\label{eq:bd26-edge-coefficients}
 A(g)=\int_0^1g(u)\,du,\qquad
 r_f=\pp_t\frac{f(1)-f(0)}t\mathbf1,\qquad
 q_g=\pp_t\bigl(tp(t)A(g)\bigr).
\end{equation}
Polynomial integration uses characteristic zero.
It is linear over $R$.
Each principal part is finite, including when $A(g)$ has arbitrarily high finite normal pole order.
The coefficients of its normal expansion are polynomial in both block positions.

Let $R_L(q)$ be the outer chiral operation on $q\boxtimes J$, retaining the first transfer module.
Let $R_R(q)$ be the corresponding operation on $J\boxtimes q$.
Both have values in $P_3(V)$.
They include transverse localization before the outer collision, followed by its principal-part map.
In particular, all normal derivatives of $q$ remain present.
For the scalar primitive,
\[
 R_L(\epsilon q)=\epsilon R_L(q),\qquad
 R_R(\epsilon q)=\epsilon R_R(q).
\]
The second equality here uses that the other state is the even current $J$.

Define a degree-zero $R$-linear map $B:O\to E$ by
\begin{equation}\label{eq:bd26-leading-map}
 \begin{aligned}
 B(fe_{ijk})&=f(ijk)[J_i|J_j|J_k] && (f\in C^0),\\
 B(cm_{ij|k})&=[\epsilon r_f|J_k] && (c|_{\{ijk,jik\}}=f),\\
 B(cm_{ij|k})&=-[q_g|J_k] && (c|_{\{ijk,jik\}}=gdu),\\
 B(cn_{k|ij})&=[J_k|\epsilon r_f] && (c|_{\{kij,kji\}}=f),\\
 B(cn_{k|ij})&=+[J_k|q_g] && (c|_{\{kij,kji\}}=gdu).
 \end{aligned}
\end{equation}
Here $f(ijk)$ is evaluation at the vertex indexed by the order $ijk$.
Set $B=0$ on positive-degree coefficients of every $e$-word and on mixed coefficients of degree at least two.
The intermediate maps in \eqref{eq:bd26-leading-map} are the binary edge assignments tensored with the remaining suspended current.
The positive sign in the last line comes from moving $du$ past that suspended current.
The unmerged vertex evaluation is specified separately and uses no endpoint identity between different edges.

\section{The full ternary cochain map}

The defect of $B$ is the degree-one map
\[
 D=QB-Bb:O\longrightarrow E.
\]
The identity $QD+Db=0$ follows from $Q^2=b^2=0$.
Every component of $D$ is supported on a proper collision diagonal.
Indeed, the projected separated map has zero internal differential and kills the coefficient and multiplication boundaries.

\begin{theorem}\label{thm:bd26-full-comparison}
Over $S_*$, the formula
\begin{equation}\label{eq:bd26-full-map}
 F=B-hD=B+\kappa^{-1}\mathsf e(QB-Bb)
 \colon O\otimes_S S_*\longrightarrow E\otimes_S S_*
\end{equation}
defines a degree-zero $R[\kappa^{-1}]$-linear cochain map.
Its separated projection is exactly the vertex insertion in \eqref{eq:bd26-leading-map}.
Its mixed intermediate components are exactly those in \eqref{eq:bd26-leading-map}.
For the mixed one-form components, its additional small-diagonal terms are
\begin{equation}\label{eq:bd26-mixed-correction}
 \begin{aligned}
 (F-B)(cm_{ij|k})&=-\kappa^{-1}[\epsilon R_L(q_g)],\\
 (F-B)(cn_{k|ij})&=+\kappa^{-1}[\epsilon R_R(q_g)].
 \end{aligned}
\end{equation}
The correction vanishes on mixed degree-zero coefficients and on mixed coefficients of degree at least two.
On unmerged words it is supported on the three proper pair diagonals and vanishes on positive-degree coefficients.
All formulas retain every normal derivative and use finite sums.
\end{theorem}
\begin{proof}
Since $D$ is closed in the Hom complex, \eqref{eq:bd26-contraction} gives
\[
 \begin{aligned}
 QF-Fb
 &=D-QhD+hDb\\
 &=D-(\mathrm{id}-hQ)D+hDb
 =h(QD+Db)=0.
 \end{aligned}
\]
The support assertion for $D$ and partition preservation by $h$ prove the separated-projection statement.

For a left mixed zero-form, endpoint regularity gives
\begin{equation}\label{eq:bd26-endpoint}
 q_{\partial_uf}=\kappa r_f.
\end{equation}
Indeed, $A(\partial_uf)=f(1)-f(0)$ and
$tp(t)-\kappa/t$ is regular and vanishes at $t=0$.
The internal intermediate differential therefore cancels $Bb$.
The remaining defect is $-[\epsilon R_L(r_f)]$ on the small diagonal.
Its image under $h$ is zero because $\epsilon^2=0$.
For a mixed one-form, the defect is $-[R_L(q_g)]$.
Applying $-h$ gives the first line of \eqref{eq:bd26-mixed-correction}.

For a right mixed zero-form, the preceding suspended current changes the internal sign.
The intermediate differential is $+\kappa[J|r_f]$ and equals $Bb$ by \eqref{eq:bd26-endpoint}.
The remaining defect is $+[\epsilon R_R(r_f)]$, which $h$ kills.
For a right mixed one-form, the defect is $+[R_R(q_g)]$.
This gives the second line of \eqref{eq:bd26-mixed-correction}.
Higher forms restrict to zero on every edge, including after their differential.

For $f\in C^0$, the defect on an unmerged word is explicitly
\begin{equation}\label{eq:bd26-unmerged-defect}
 D(fe_{ijk})=
 Q_2\bigl(f(ijk)[J_i|J_j|J_k]\bigr)
 +B(f\omega_{ij}m_{ij|k})+B(f\omega_{jk}n_{i|jk}).
\end{equation}
All three terms lie on proper pair diagonals.
The coefficient differential gives a positive-degree unmerged coefficient, which $B$ kills.
If $c$ has positive degree, both $B(ce_{ijk})$ and $B(d_Cc\,e_{ijk})$ vanish.
The other boundary coefficients $c\omega$ have degree at least two and vanish on edges.
Hence $D(ce_{ijk})=0$ in these degrees.
Formula \eqref{eq:bd26-unmerged-defect} and the explicit chiral maps make \eqref{eq:bd26-full-map} a finite formula on every input.
\end{proof}

The general Hom identity also gives an explicit null homotopy of $F$:
\begin{equation}\label{eq:bd26-null}
 K=hF=hB,
 \qquad QK+Kb=F.
\end{equation}
Here $h^2=0$ gives the first equality and $QF=Fb$ gives the second.
This homotopy is consistent with the contractibility of the nonzero-level target.

For $e=e_{123}$ with constant coefficient, the two adjacent terms in \eqref{eq:bd26-unmerged-defect} cancel.
Their coefficients are $\pp_t p(t)=\kappa t^{-2}\mathbf1$.
The remaining defect is the nonadjacent collision:
\begin{equation}\label{eq:bd26-constant-word}
 D(e)=-\kappa[\mathbf1\delta_{13}\partial_{t_{13}}|J_2],
 \qquad
 F(e)=[J_1|J_2|J_3]-[\epsilon\delta_{13}\partial_{t_{13}}|J_2],
 \quad t_{13}=z_1-z_3.
\end{equation}
The second formula is polynomial in the level.
Its partial-diagonal internal differential cancels the $13$ component.
Its next collision vanishes because the vacuum-current outer coefficient is regular.
The correction is supported on $\Delta_{13}$, including its locus away from both adjacent diagonals.
Thus the cancellation occurs on the support of the actual nonadjacent defect.

For each label $i$, let $\mathsf e_i$ insert $\epsilon$ in the block containing $i$, with the sign of \eqref{eq:bd26-epsilon-operator}.
The operators satisfy $\mathsf e_i\mathsf e_j+\mathsf e_j\mathsf e_i=0$.
For distinct blocks this is the tensor sign for two odd maps.
For the same block it follows from $\epsilon^2=0$.
Each operator also satisfies $Q\mathsf e_i+\mathsf e_iQ=-\kappa$.

\begin{corollary}\label{cor:bd26-equivariant}
Replace $h$ in \eqref{eq:bd26-full-map} by
\[
 h_{\mathrm{sym}}=-\frac1{3\kappa}(\mathsf e_1+\mathsf e_2+\mathsf e_3).
\]
The resulting map $F_{\mathrm{sym}}$ is equivariant under permutation of the three labels, including the coordinate and suspended-state signs.
It has the same mixed components as $F$.
The maps satisfy
\[
 F_{\mathrm{sym}}-F=QL+Lb,
 \qquad L=(h_{\mathrm{sym}}-h)B.
\]
\end{corollary}
\begin{proof}
The preceding identities give $h_{\mathrm{sym}}^2=0$ and
$Qh_{\mathrm{sym}}+h_{\mathrm{sym}}Q=\mathrm{id}$.
The proof of Theorem~\ref{thm:bd26-full-comparison} therefore applies.
The edge formulas, vertex evaluation, and chiral maps in $B$ commute with label permutations.
Permutations exchange the three operators $\mathsf e_i$, so their average also commutes with them.
On the small diagonal, all three scalar operators coincide, which preserves the mixed correction.
Finally, $Q(h_{\mathrm{sym}}-h)=-(h_{\mathrm{sym}}-h)Q$ gives
$QL+Lb=-(h_{\mathrm{sym}}-h)(QB-Bb)=F_{\mathrm{sym}}-F$.
\end{proof}

\section{Polynomial comparison on the global edge domain}

The global order-edge ring $R[t^{-1}]$ contains no inverse of an outer position difference.
Consequently, a mixed principal part in \eqref{eq:bd26-edge-coefficients} has polynomial coefficients in the outer variable.
This gives a divisibility statement at every inner pole order.

\begin{lemma}\label{lem:bd26-global-divisibility}
For every global mixed edge one-form,
\[
 R_L(q_g)\in\kappa P_3(V),\qquad R_R(q_g)\in\kappa P_3(V).
\]
The quotients by $\kappa$ are unique polynomial elements of the finite principal-parts module.
\end{lemma}
\begin{proof}
Expand $q_g$ in the normal basis at its pair diagonal.
Every normal coefficient is a polynomial in the two block positions with a state in $V$.
Reduction modulo $\kappa$ makes the vertex field
$Y_0(a,s)b=(e^{sT}a)b$.
Multiplication by a polynomial outer coefficient remains regular in $s$.
Its outer principal part is therefore zero.
This argument applies to every inner normal derivative, since the transfer retains these derivatives and regular coefficients remain regular under differentiation.
Hence both outer outputs vanish modulo $\kappa$.
The module $P_3(V)$ is free over $S$ in the monomial-state and normal bases.
It follows that each output has a unique quotient by $\kappa$.
\end{proof}

Thus \eqref{eq:bd26-mixed-correction} defines a polynomial comparison on the entire global mixed sector.
There is no bound on its finite inner pole order.
At level zero, the quotient is taken in $S$ before specialization.
It is not defined by numerical division at $\kappa=0$.

For an unmerged word $e_{ijk}$, consider its two adjacent edges.
Let $C^0_{ijk,\mathrm{reg}}$ consist of $f\in C^0$ whose average on each of these edges is regular in the corresponding normal variable.
Equivalently, its average belongs to $R$ inside $R[t^{-1}]$ on that edge.
Define an $R$-subcomplex $O_{\mathrm{reg}}\subset O$ by using these zero-form coefficients for the unmerged words.
Keep every positive-degree unmerged coefficient and every mixed coefficient.
It is a subcomplex by \eqref{eq:bd26-source}.
It is not asserted to be a $C$-submodule.

\begin{theorem}\label{thm:bd26-polynomial}
Formula \eqref{eq:bd26-full-map} restricts to an $S$-polynomial cochain map
\[
 F_{\mathrm{reg}}:O_{\mathrm{reg}}\longrightarrow E.
\]
It specializes to every numerical level, including zero.
Its mixed sector contains every global inner pole order, and its unmerged sector contains every constant-coefficient ordered word.
\end{theorem}
\begin{proof}
Lemma~\ref{lem:bd26-global-divisibility} proves polynomiality on mixed one-forms.
The mixed zero-form correction vanishes.
For an unmerged zero-form, all pair collisions of the regular vertex coefficient $f(ijk)$ are divisible by $\kappa$.
Indeed, modulo $\kappa$ the regular-input vertex product has no pole.
On an adjacent edge, the coefficient of $f\omega_{ij}$ is $f(u)du/t$.
Its averaged inner output is
\[
 q_{f/t}=\pp_t\bigl(p(t)A(f)\bigr).
\]
The regular-average condition makes its regular field contribution vanish under $\pp_t$.
Thus it is also divisible by $\kappa$.
All components of \eqref{eq:bd26-unmerged-defect} are consequently divisible by $\kappa$.
This proves that $\mathsf e D/\kappa$ is polynomial on $O_{\mathrm{reg}}$.
The correction vanishes on positive-degree unmerged coefficients.

The cochain identity holds after inverting $\kappa$ by Theorem~\ref{thm:bd26-full-comparison}.
Every summand of $E$ is $S$-torsion-free: localization and the finite normal forms preserve this property.
The identity therefore holds over $S$.
Tensoring this polynomial identity with $S/(\kappa-k)$ proves every stated specialization.
\end{proof}

The same polynomial statement holds for $F_{\mathrm{sym}}$.
The proof uses only divisibility of $D$ on unmerged zero-forms and divisibility of the small-diagonal mixed defect.
It therefore applies to each $\mathsf e_i$ and to their average.

There is a global coefficient outside the regular-average subcomplex:
\[
 f=\frac{q_{12}(1-q_{12})}{t},\qquad t=z_1-z_2.
\]
When the relative order of $1,2$ is fixed, its numerator is zero.
On every other component, $t$ is inverted.
Thus $f$ is an actual compatible degree-zero family.
On the edge $\{123,213\}$, its average is $1/(6t)$.
On the edge $\{123,132\}$, its restriction is zero.
Its values at all six vertices are zero.

\begin{proposition}\label{prop:bd26-positive-form-obstruction}
No $S$-polynomial degree-zero cochain map on all of $O$ can simultaneously retain the separated and mixed intermediate components of $B$ and kill all positive-degree unmerged coefficients.
The same statement holds at level zero.
\end{proposition}
\begin{proof}
Use $fe_{123}$ with the displayed global family.
Its source boundary is
\[
 b(fe_{123})=(d_Cf)e_{123}
       -f\omega_{12}m_{12|3}-f\omega_{23}n_{1|23}.
\]
The first term has zero image under the stated hypothesis.
The last term has zero mixed intermediate image because $f$ vanishes on its edge.
For the middle term,
\[
 q_{f/t}=\pp_t\left(\frac{p(t)}{6t}\right)
 =\frac\kappa{6t^3}\mathbf1+\frac{J^2}{6t}.
\]
Thus the partial-diagonal component of the boundary image reduces to
$[J^2\delta_{12}/6|J_3]\ne0$ modulo $\kappa$.
The proposed image of $fe_{123}$ has zero separated component, since all its vertex values vanish.
Its degree is $-3$.
In this degree, the proper pair diagonals contain one $\epsilon$ factor, while the small-diagonal summand is zero.
The partial-diagonal component of its differential is consequently divisible by $\kappa$.
This contradicts the displayed nonzero reduction of the source-boundary image.
At level zero that target differential is zero, giving the same contradiction.
\end{proof}

A polynomial extension to every unmerged coefficient with these intermediate images must therefore include nonzero positive-degree unmerged components.
This condition concerns the map on the original ordered complex, before differential induction.

The coefficient-domain hypothesis is essential.
On a mixed domain admitting the outer coefficient $s^{-1}$, take
$gdu=du/(t^2s)$ for the left placement.
Then
\[
 q_g=\frac\kappa{st^3}\mathbf1+\frac{J^2}{st}.
\]
The actual outer output is
\begin{equation}\label{eq:bd26-outer-obstruction}
 R_L(q_g)=
 \frac{\kappa J}{t^3s}+\frac{J^3}{ts}
 +\frac{2\kappa J}{ts^3}
 +\frac{2\kappa x_2}{ts^2}
 +\frac{2\kappa x_3}{ts}.
\end{equation}
To obtain the last four terms, use
$Y(J^2,s)J=(e^{sT}J)^2J+2\kappa s^{-2}e^{sT}J$.
The first term retains the inner derivative $\partial_t^2/2$ through the unit operation.
Modulo $\kappa$, \eqref{eq:bd26-outer-obstruction} is $J^3/(ts)\ne0$.
Thus the polynomial mixed formula does not extend to this outer-pole domain while retaining its intermediate image and allowing only a small-diagonal correction.
On the closed degree-$-2$ input, the target correction has the form $[\epsilon v]$ and differential $-\kappa[v]$.
This cannot cancel the nondivisible coefficient.
The global divisibility lemma and this obstruction have different coefficient domains.

\section{Differential induction and the remaining comparison}

The map $F$ is linear over position functions.
The ordered symbols in $O$ have not been assigned a connection commuting with their multiplication differential.
A precise right differential module source can instead be constructed by induction.
Let $\mathcal D_R$ be the Weyl algebra over $S$, generated by $R$ and $\partial_{z_1},\partial_{z_2},\partial_{z_3}$.
Its relations are $[\partial_{z_i},z_j]=\delta_{ij}$ and $[\partial_{z_i},\partial_{z_j}]=0$.
All Weyl generators have degree zero.
For an $R$-complex $M$, put
\begin{equation}\label{eq:bd26-induction}
 \operatorname{Ind}(M)=M\otimes_R\mathcal D_R,
 \qquad d(m\otimes P)=dm\otimes P.
\end{equation}
Its right action is multiplication on the second factor.
Normal ordering expresses each element uniquely as a finite sum
$\sum_\alpha m_\alpha\otimes\partial^\alpha$.
Thus $\mathcal D_R$ is free as a left $R$-module and induction is exact.
The source contains the specified formal derivatives of every state-coefficient expression.

\begin{proposition}\label{prop:bd26-induced-comparison}
There are cochain maps of right differential modules
\[
 \begin{aligned}
 F^{\mathcal D}:\operatorname{Ind}(O)\otimes_S S_*&\longrightarrow E\otimes_S S_*,\\
 F_{\mathrm{reg}}^{\mathcal D}:\operatorname{Ind}(O_{\mathrm{reg}})&\longrightarrow E,
 \end{aligned}
 \qquad m\otimes P\longmapsto F(m)P.
\]
They retain every displayed normal-derivative identity.
The first has right differential module null homotopy
$K^{\mathcal D}(m\otimes P)=hB(m)P$.
\end{proposition}
\begin{proof}
Function-linearity gives $F(mr)P=F(m)rP$, so the formula is balanced over $R$.
Right differential linearity follows from associativity of the right action.
The target differential commutes with that action.
Hence $Q(F(m)P)=F(bm)P$, proving the cochain equation.
The same argument applies to the polynomial subcomplex.
Apply \eqref{eq:bd26-null} and the right differential linearity of $h$ to prove the homotopy formula.
\end{proof}

Induction specifies a new source rather than a connection on the original finite list of ordered symbols.
It is not an identification of the underlying $R$-complexes $O$ and $\operatorname{Ind}(O)$.
For example, their formal derivative summands are separate in the displayed normal form.
A comparison on a quotient imposing state-translation relations must show that $F^{\mathcal D}$ kills those relations.

The state complex $W$ has $H^{-1}(W)=0$ and $H^0(W)=V/\kappa V$ over $S$.
At nonzero numerical level, multiplication by $\epsilon/\kappa$ contracts it.
At zero level, its ordinary specialization has zero differential and state space $V_0\oplus\epsilon V_0$.
Therefore neither \eqref{eq:bd26-full-map} nor differential induction gives a quasi-isomorphic replacement of the nonzero-level Heisenberg states.
The ternary map also specifies no compatibility with deconcatenation, ordered surjections, or four-input compositions.
Such compatibility requires new equations on those maps and their actual coefficient domains.
The partial-diagonal term \eqref{eq:bd26-constant-word} supplies the missing three-input chain equation with the stated mixed images.

\chapter{Polynomial comparison for every global ternary coefficient}
\label{chap:bd26-polynomial-full}

\section{The weighted edge integral}

The family $q_{12}(1-q_{12})/(z_1-z_2)$ forces a nonzero image of an unmerged one-form.
The weighted edge integral supplies that image.
It cancels the regular composite states in the adjacent comparison defect before any division by the level.

Retain the base $S$, the ordered complex $O$, the five-partition complex $(E,Q)$, and the chiral operations of Chapter~\ref{chap:bd26-comparison}.
In particular, the source coefficients are global compatible families.
Each adjacent edge ring is $R[t^{-1}]$ and has no outer pole.
All principal parts and all differential-operator sums remain finite.

On an oriented edge with coordinate $u=q_{ij}$, define
\begin{equation}\label{eq:bd26-weighted-integral}
 M(g)=\int_0^1u g(u)\,du,
 \qquad v_g=\pp_t\bigl(p(t)M(g)\bigr),
 \qquad t=z_i-z_j.
\end{equation}
The degree of $M(g)$ is zero after the coefficient one-form $gdu$ has been integrated.
Every normal coefficient of $v_g$ is polynomial in the two remaining block positions.
Polynomial integration by parts gives
\begin{equation}\label{eq:bd26-weighted-endpoint}
 M(\partial_uf)=f(1)-A(f),\qquad
 v_{\partial_uf}+q_{f/t}=\pp_t\bigl(p(t)f(1)\bigr).
\end{equation}
The second identity holds at every finite inner pole order.
It is an identity before any outer collision is taken.

For $c\in C^1$ and the word $e_{ijk}$, write its two actual edge restrictions as
\[
 c|_{\{ijk,jik\}}=g_Ldu_L,
 \qquad c|_{\{ijk,ikj\}}=g_Rdu_R,
 \qquad u_L=q_{ij},\quad u_R=q_{jk}.
\]
Both edges have $u=1$ at the vertex $ijk$.
Let $v_L$ and $v_R$ be their principal parts from \eqref{eq:bd26-weighted-integral}, in their respective normal coordinates.
Define a degree-zero $R$-linear map $B^+:O\to E$ by retaining every mixed component of $B$ and setting
\begin{equation}\label{eq:bd26-positive-leading}
 \begin{aligned}
 B^+(fe_{ijk})&=f(ijk)[J_i|J_j|J_k] && (f\in C^0),\\
 B^+(ce_{ijk})&=[v_L|J_k]-[J_i|v_R] && (c\in C^1),\\
 B^+(ce_{ijk})&=0 && (|c|\geq2).
 \end{aligned}
\end{equation}
The first mixed block in the second line lies on $\Delta_{ij}$.
The second lies on $\Delta_{jk}$.
Each has total degree $-2$, equal to the degree of $ce_{ijk}$.
The minus sign in the right term is the sign of the adjacent $jk$ collision when the remaining suspended current is placed first.

\section{Divisibility of the complete comparison defect}

The defect of $B^+$ is
\[
 D^+=QB^+-B^+b.
\]
It satisfies $QD^++D^+b=0$.
For an outer operation, let a finite inner principal part be called regular in the outer variable when its normal coefficients are polynomials in the block positions.
This property concerns its coefficients before the outer chiral operation.
The outer output itself can have poles.

\begin{lemma}\label{lem:bd26-arbitrary-regular-outer}
For any finite inner principal part $v$ with states in $V$ that is regular in the outer variable,
\[
 R_L(v)\in\kappa P_3(V),\qquad R_R(v)\in\kappa P_3(V).
\]
\end{lemma}
\begin{proof}
Reduce its monomial-state and normal coefficients modulo $\kappa$.
The outer fields are $Y_0(a,s)b=(e^{sT}a)b$.
The coefficient is polynomial in $s$, so the outer principal part is zero.
Every retained inner normal derivative preserves this conclusion, since it differentiates regular coefficients or remains a normal derivative in the final transfer module.
The output therefore vanishes modulo $\kappa$.
The free basis \eqref{eq:bd26-small} gives the asserted divisibility.
\end{proof}

\begin{proposition}\label{prop:bd26-complete-defect}
Every value of $D^+$ belongs to $\kappa E$.
Its values on unmerged words are
\begin{equation}\label{eq:bd26-positive-defect}
 \begin{aligned}
 D^+(fe_{ijk})&=-[\mu_{ik}(f(ijk)(J_i,J_k))|J_j],
                    &&f\in C^0,\\
 D^+(ce_{ijk})&=[R_L(v_L)-R_R(v_R)],
                    &&c\in C^1,\\
 D^+(ce_{ijk})&=0, &&|c|\geq2.
 \end{aligned}
\end{equation}
In the first line, the coefficient $f(ijk)$ is retained as a function of all three positions during the first collision.
The first output lies on $\Delta_{ik}$.
The second output lies on the small diagonal, with both terms written using the same density and normal-coordinate identification.
\end{proposition}
\begin{proof}
For a zero-form $f$, the $ij$ component of $QB^+(fe_{ijk})$ is
$[\pp_t(p(t)f(ijk))|J_k]$.
The corresponding component of $B^+b(fe_{ijk})$ is
\[
 [v_{\partial_uf}+q_{f/t}|J_k].
\]
The value of $f$ at $u=1$ on this edge is $f(ijk)$.
Equation~\eqref{eq:bd26-weighted-endpoint} makes the two components equal.

For the $jk$ component, both expressions have a minus sign with $J_i$ in the first block.
The same endpoint identity proves their equality.
The source differential has no $ik$ face, so the nonadjacent chiral term remains.
The three-current unshuffle sign for that term is negative.
This proves the first line of \eqref{eq:bd26-positive-defect}.

The coefficient $f(ijk)$ is regular in every normal variable.
Thus \eqref{eq:bd26-field} gives
\[
 \mu_{ik}(f(ijk)(J_i,J_k))
 =\kappa\pp_{t_{ik}}\frac{f(ijk)}{t_{ik}^2}\mathbf1.
\]
This proves divisibility on every unmerged zero-form.

For $c\in C^1$, the states in $B^+(ce_{ijk})$ are even and have zero internal differential.
Their collision differential is $[R_L(v_L)-R_R(v_R)]$.
The coefficient differential $d_Cc$ has degree two, so its unmerged image is zero.
The other boundary coefficients are $c\omega$, whose degree-two edge restrictions are zero.
Hence $B^+b(ce_{ijk})=0$.
This proves the second line of \eqref{eq:bd26-positive-defect}.
Lemma~\ref{lem:bd26-arbitrary-regular-outer} proves divisibility of both terms.
All higher-degree unmerged inputs have zero defect by the same edge-degree argument.

The mixed values of $D^+$ equal those of $D$.
On mixed zero-forms, their possible outer term contains $r_f$, a vacuum principal part with polynomial outer coefficients.
The vacuum operation with that outer coefficient has no pole, so this term is zero.
On mixed one-forms, the values are $-[R_L(q_g)]$ and $+[R_R(q_g)]$.
They are divisible by $\kappa$ by Lemma~\ref{lem:bd26-global-divisibility}.
Higher-degree mixed restrictions vanish.
This checks every summand of $O$.
\end{proof}

\section{The polynomial ternary map}

Every summand of $E$ is torsion-free over $S$.
Proposition~\ref{prop:bd26-complete-defect} therefore defines a unique degree-one $R$-linear map
\[
 T^+:O\longrightarrow E,
 \qquad \kappa T^+=D^+.
\]
Its values are obtained by division of polynomial coefficients known to be divisible by $\kappa$.
This definition uses no inverse of $\kappa$ in the coefficient ring.

\begin{theorem}\label{thm:bd26-polynomial-complete}
The formula
\begin{equation}\label{eq:bd26-polynomial-complete}
 G=B^++\mathsf eT^+:O\longrightarrow E
\end{equation}
defines a degree-zero $R$-linear cochain map over $S$.
It is defined on all six unmerged words and all twelve mixed words, with every global compatible-family coefficient and every finite inner pole order.
Its mixed components are exactly those of \eqref{eq:bd26-leading-map} and \eqref{eq:bd26-mixed-correction}, with the polynomial quotients taken before specialization.
Its separated projection is $fe_{ijk}\mapsto f(ijk)[J_i|J_j|J_k]$.
Its positive-degree unmerged components are determined by \eqref{eq:bd26-positive-leading} and \eqref{eq:bd26-positive-defect}.
It specializes to every numerical level, including zero.
\end{theorem}
\begin{proof}
The closed-defect equation and $S$-torsion-freeness imply
\[
 QT^++T^+b=0.
\]
Using $Q\mathsf e+\mathsf eQ=-\kappa$ gives
\[
 \begin{aligned}
 QG-Gb
 &=D^++Q\mathsf eT^+-\mathsf eT^+b\\
 &=\kappa T^+-\kappa T^+-\mathsf e(QT^++T^+b)=0.
 \end{aligned}
\]
This proves the chain equation over $S$ directly.
The complete list of defect components proves the support and component assertions.
Every scalar quotient, integration, and residue is $R$-linear.
All outputs have finite normal expansions.
Finally, specialize the polynomial identity $QG=Gb$ along $S\to S/(\kappa-k)$.
No division at the numerical value $k=0$ occurs.
\end{proof}

In particular, the unmerged one-form image is
\begin{equation}\label{eq:bd26-polynomial-oneform}
 G(ce_{ijk})=[v_L|J_k]-[J_i|v_R]
   +\left[\epsilon\,\frac{R_L(v_L)-R_R(v_R)}\kappa\right].
\end{equation}
The displayed quotient is the unique polynomial quotient.
For a constant unmerged word, $G$ has the same nonadjacent correction as \eqref{eq:bd26-constant-word}.

The necessity of the positive-degree component has a direct check.
Take $f=q_{12}(1-q_{12})/t$ and $e=e_{123}$.
Then $G(fe)=0$ because every vertex value vanishes.
On the left edge,
\[
 M(\partial_uf)=-\frac1{6t},\qquad
 v_{\partial_uf}=-\frac\kappa{6t^3}\mathbf1-\frac{J^2}{6t}
 =-q_{f/t}.
\]
The right edge restriction is zero.
Therefore the proper-diagonal terms of
$G((d_Cf)e)-G(f\omega_{12}m_{12|3})$ cancel.
Their small-diagonal terms also cancel, by the same identity after the outer operation.
Thus $G(b(fe))=0=QG(fe)$, including the nonzero composite state modulo $\kappa$.

\begin{corollary}\label{cor:bd26-polynomial-equivariant}
Replacing $\mathsf e$ in \eqref{eq:bd26-polynomial-complete} by
$(\mathsf e_1+\mathsf e_2+\mathsf e_3)/3$ gives a permutation-equivariant polynomial cochain map $G_{\mathrm{sym}}$.
There are also polynomial cochain maps of right differential modules
\[
 \operatorname{Ind}(O)\longrightarrow E,
 \qquad m\otimes P\longmapsto G(m)P
 \quad\hbox{or}\quad G_{\mathrm{sym}}(m)P.
\]
\end{corollary}
\begin{proof}
Each scalar insertion has anticommutator $-\kappa$ with $Q$.
Their average has the same identity and commutes with label permutations.
The two oriented edge integrals in $B^+$ commute with relabeling, since each has its ordered-word vertex at $u=1$.
The defect and its unique polynomial quotient are therefore equivariant.
This proves the first assertion.
The induction formulas are balanced over $R$ and commute with the differential by Proposition~\ref{prop:bd26-induced-comparison}.
\end{proof}

\section{A comparison with the unextended level-zero chiral algebra}

Put $V_0=\mathbb C[x_1,x_2,\ldots]$ with field
$Y_0(a,t)b=(e^{tT}a)b$ and zero state differential.
Its chiral operation on pole-valued inputs is still the principal-part operation \eqref{eq:bd26-chiral}.
Let $E_0(V_0)$ denote its five-partition Chevalley--Cousin complex, using the same shifts and densities.
Let $O_0=O\otimes_S S/(\kappa)$.

The projection $W\otimes_S S/(\kappa)\to V_0$ sends $\epsilon$ to zero and preserves every vertex operation.
Applying it to each partition factor gives a cochain map to $E_0(V_0)$.
Its composite with the specialization of $G$ is denoted $G_0$.

\begin{corollary}\label{cor:bd26-native-zero}
There is a permutation-equivariant cochain map
\[
 G_0:O_0\longrightarrow E_0(V_0)
\]
to the unextended level-zero chiral algebra.
It is obtained from $B^+$ by setting $\kappa=\epsilon=0$.
It also induces a cochain map of right differential modules from $\operatorname{Ind}(O_0)$.
The degree-$-3$ source cycle
\begin{equation}\label{eq:bd26-nonzero-cycle}
 z=(z_1-z_2)(z_2-z_3)e_{123}
   +(z_2-z_3)q_{12}m_{12|3}
   +(z_1-z_2)q_{23}n_{1|23}
\end{equation}
has image representing a nonzero class in $H^{-3}(E_0(V_0))$.
\end{corollary}
\begin{proof}
The state projection is a differential graded vertex map because both differentials are zero and the scalar factor $\epsilon$ is an ideal.
It therefore commutes with the chiral operation, its transfers, and the Chevalley--Cousin differential.
All correction terms in $G-B^+$ contain $\epsilon$ and vanish under the projection.
Hence $G_0$ is exactly the displayed specialization of $B^+$.
The same specialization of the equivariant map $G_{\mathrm{sym}}$ gives $G_0$, proving equivariance.
Differential induction applies as before.

Equation~\eqref{eq:bd26-source} and $d_Cq_{ij}=(z_i-z_j)\omega_{ij}$ give $bz=0$.
The mixed zero-form images contain $\epsilon$ and therefore vanish under the projection.
Consequently,
\[
 G_0(z)=(z_1-z_2)(z_2-z_3)[J_1|J_2|J_3].
\]
This is nonzero in the separated summand.
It is closed because every regular-input level-zero chiral bracket vanishes.
The complex $E_0(V_0)$ has terms only in degrees $-3,-2,-1$.
There is no incoming boundary in degree $-3$.
Thus the displayed image represents a nonzero cohomology class.
\end{proof}

This comparison uses the unextended state space only at level zero.
It asserts no quasi-isomorphism of the full complexes.

The source of this polynomial map is the global ordered coefficient complex.
Admitting independent outer poles changes the comparison problem.
The cubic obstruction \eqref{eq:bd26-outer-obstruction} still excludes the specified extension on that larger mixed domain with its intermediate images fixed.
The polynomial family before the level-zero projection uses the Koszul state extension $W$.
It preserves neither the nonzero-level diagonal state cohomology of $V$ nor a presently unspecified system of state-translation relations on $O$.
Compatibility with ordered surjections, deconcatenation, and four-input operations remains an additional requirement.

\begin{thebibliography}{9}
\bibitem{BDdiagonal}
A. Beilinson and V. Drinfeld,
\emph{Chiral Algebras}, American Mathematical Society, 2004,
\S3.4.11, equation (3.4.11.1), pp.~182--183.
\end{thebibliography}
\end{document}
